% GENERATED FILE. Edit canonical lessons, then run npm run generate:books.
% canonical-source-hash: fnv1a64:1f67f2a497338e5b
% canonical-lessons: ZH-C148-xia-ge-yue, ZH-W148-yue-month, ZH-C148-yao-future, ZH-C148-xia-xingqi, ZH-C148-shang-ge-yue, ZH-C148-dasuan, ZH-C148-jianglai, ZH-C148-practice

\chapter{Next Month}
\label{ch:zh-next-month}
\hlchaptermodality{148}

\begin{chapteropening}
\textbf{By the end of this chapter:} \emph{I can say what I am going to do next week and next month.}

Saying what you are going to do with \zh{要} and dǎsuàn, next week, next month and last month, and what is about to happen with \zh{要}…\zh{了}, built on one new written character: \zh{月}.
\end{chapteropening}

\section[\texorpdfstring{\zh{下个月} (xià ge yuè)}{下个月 (xià ge yuè)}]{\zh{下个月} (xià ge yuè) — next month}

\label{lesson:ZH-C148-xia-ge-yue}



\begin{quote}
A new chapter, and it opens a little further ahead than tomorrow.

\begin{itemize}
\raggedright
  \item \emph{Recall:} say \zh{明天会下雨吗} — \emph{míngtiān huì xiàyǔ ma}, and answer it — \textbf{R1}, the chapter you have finished
  \item \emph{Recall:} say \emph{tiānqì yùbào}, the weather forecast — \textbf{R1}
  \item \emph{Recall:} say \emph{yuè}, a month — \textbf{R3}, durable
  \item \emph{Recall:} say \zh{明年} — \emph{míngnián}, next year — \textbf{R3}
\end{itemize}
\end{quote}

\subsection*{You'll want to know: \zh{下个月}}
\textbf{\zh{下个月}} — \emph{xià ge yuè} — \textquotedblleft{}next month\textquotedblright{}.

\begin{quote}
\textbf{\zh{下个月我们去}} \emph{xià ge yuè wǒmen qù} — We're going next month.
\end{quote}

The next lesson puts the last character in your hand. It is the \emph{yuè} you already say for a month.

\begin{cousinweb}
\textbf{\zh{下个月}} (\emph{xià ge yuè}) is built from \zh{下} 'below; after', \zh{个}, the measure word, and \zh{月} 'month'. All three are inherited from Middle Chinese, \emph{haeX}, \emph{kaH} and \emph{ngjwot} (Guangyun, in Baxter's transcription). The Guangyun glosses \emph{haeX} with, among other things, 'after' and 'the bottom', and quotes an old saying about the moon under \emph{ngjwot}; 'next' in a row of months is assumed from 'after', not quoted.
\end{cousinweb}

\begin{grammarlens}[title={What you've built}]
\zh{下} here is the month after this one. Like \zh{明天} and \zh{后天}, \zh{下个月} goes early in the sentence: before the verb, or before the person.

\zh{明年} says next year with \zh{明}; the month and the week take \zh{下}.
\end{grammarlens}

\subsection*{Your turn}
\begin{itemize}
\raggedright
  \item \emph{Say it:} \emph{xià ge yuè}
  \item \emph{Say it:} \emph{xià ge yuè}, once more
  \item \emph{Say it:} \zh{下个月我们去} — \emph{xià ge yuè wǒmen qù} — We're going next month.
  \item \emph{Say it:} \zh{明年} — \emph{míngnián} — next year
  \item \emph{Recall:} say \emph{yuè}, then \emph{xià ge yuè}
\end{itemize}

\begin{tcolorbox}[breakable,colback=teal!4,colframe=teal!35!black,title={Before you move on}]
What does \zh{下个月} mean? (\textquotedblleft{}Next month\textquotedblright{}.) And next year? (\textbf{\zh{明年}}.)
\end{tcolorbox}
\section[\texorpdfstring{\zh{月} (yuè)}{月 (yuè)}]{\zh{月} — the moon piece, standing on its own}

\label{lesson:ZH-W148-yue-month}



\begin{quote}
The word first, then the character inside it.

\begin{itemize}
\raggedright
  \item \emph{Recall:} say \emph{xià ge yuè} — \textbf{R1}, one lesson back, before it decays
\end{itemize}

\begin{quote}
\zh{明} \zh{期} \zh{前} \zh{能}
\end{quote}

\emph{Read it:} each one aloud, with nothing said first — \emph{míng}, \emph{qī}, \emph{qián}, \emph{néng}
\end{quote}

\subsection*{Script — the moon piece, standing on its own}
\hlblockfigure{\detokenize{figures/ZH-W148-yue-month-filmstrip.pdf}}{How \zh{月} is written, stroke by stroke}

\begin{quote}
\zh{月}
\end{quote}

\textbf{Nothing in it is new.} You have written it four times already: on the right of \zh{明} and of \zh{期}, under the top of \zh{前}, and on the left of \zh{能}. Now it stands alone, a little taller.

\textbf{A left-falling stroke} down the left side, curving out to the lower left.

\textbf{The hooked frame}: a horizontal across the top that turns down the right side and hooks to the left at its foot.

\textbf{Two inner bars}, the upper one first, each from the left side across to the right.

1 + 1 + 2 = four.

\subsection*{Your turn}
\begin{itemize}
\raggedright
  \item \emph{Write it:} \zh{月}, all four, with the model in front of you
  \item \emph{Cover:} write it again from memory, and count the pieces, not the strokes
\end{itemize}

\begin{tcolorbox}[breakable,colback=teal!4,colframe=teal!35!black,title={Before you move on}]
How many strokes in \zh{月}? (\textbf{Four}.) Where have you written it before? (\textbf{In \zh{明}, \zh{期}, \zh{前} and \zh{能}}.) Which stroke goes in last? (\textbf{The lower inner bar}.)

Source: \href{https://raw.githubusercontent.com/chanind/hanzi-writer-data/68d10a4b21150cae5e1ebbd223eed289cf32d90c/data/\%E6\%9C\%88.json}{Hanzi Writer Data \zh{月}, pinned stroke paths and medians}.
\end{tcolorbox}
\section[\texorpdfstring{\zh{火车要开了} (huǒchē yào kāi le)}{火车要开了 (huǒchē yào kāi le)}]{\zh{火车要开了} — \zh{要} for what is coming}

\label{lesson:ZH-C148-yao-future}



\begin{quote}
A word you know, before a verb this time.

\begin{itemize}
\raggedright
  \item \emph{Read it:} \zh{下个月}, aloud — \emph{xià ge yuè}
  \item \emph{Recall:} say \zh{我要水} — \emph{wǒ yào shuǐ}, I want some water
  \item \emph{Recall:} say \zh{火车} — \emph{huǒchē}, a train
\end{itemize}
\end{quote}

\begin{grammarlens}[title={Grammar Lens: \zh{要} before a verb, for what is coming}]
\begin{quote}
\textbf{\zh{要} before a thing says what you want. \zh{要} before a verb says what you are going to do, and with \zh{了} at the end, what is about to happen.}
\end{quote}

\begin{quote}
\textbf{\zh{我要水}} \emph{wǒ yào shuǐ} — I want some water. \textbf{\zh{下个月我要去中国}} \emph{xià ge yuè wǒ yào qù Zhōngguó} — Next month I'm going to China.
\end{quote}

The test is the word after \zh{要}: a thing, and you want it; a verb, and you are going to do it.

\begin{quote}
\textbf{\zh{火车要开了}} \emph{huǒchē yào kāi le} — The train is about to leave.
\end{quote}

The \zh{了} at the end is the \zh{了} of a change, as in \zh{我是大学生了}: the moment is nearly here. For any second now, put \emph{kuài} in front: \emph{huǒchē kuàiyào kāi le}.
\end{grammarlens}

\subsection*{Your turn}
Say these aloud:

\begin{itemize}
\raggedright
  \item \zh{火车要开了} — \emph{huǒchē yào kāi le} — The train is about to leave.
  \item \zh{下个月我要去中国} — \emph{xià ge yuè wǒ yào qù Zhōngguó} — Next month I'm going to China.
  \item \zh{我要水} — \emph{wǒ yào shuǐ} — I want some water.
  \item \emph{huǒchē kuàiyào kāi le} — the train is leaving any second
\end{itemize}

\begin{tcolorbox}[breakable,colback=teal!4,colframe=teal!35!black,title={Before you move on}]
What does \zh{要} mean before a thing? (\textbf{Want}.) Before a verb? (\textbf{Going to}.) What does \zh{了} add to \zh{火车要开了}? (\textbf{It is about to happen}.)
\end{tcolorbox}
\section[\texorpdfstring{\zh{下星期} (xià xīngqī)}{下星期 (xià xīngqī)}]{\zh{下星期} (xià xīngqī) — next week}

\label{lesson:ZH-C148-xia-xingqi}



\begin{quote}
Before the new one: read the last lesson's sentence, then two words for the week.

\begin{itemize}
\raggedright
  \item \emph{Read it:} \zh{火车要开了}, aloud — \emph{huǒchē yào kāi le}
  \item \emph{Recall:} say \zh{星期天} — \emph{xīngqītiān}, Sunday — \textbf{R3}, durable
  \item \emph{Recall:} say \emph{zhōumò}, the weekend — \textbf{R3}
\end{itemize}
\end{quote}

\subsection*{You'll want to know: \zh{下星期}}
\textbf{\zh{下星期}} — \emph{xià xīngqī} — \textquotedblleft{}next week\textquotedblright{}.

\begin{quote}
\textbf{\zh{下星期我们去看书}} \emph{xià xīngqī wǒmen qù kàn shū} — Next week we're going to go and read.
\end{quote}

\begin{quote}
\textbf{\zh{下星期我要开车去}} \emph{xià xīngqī wǒ yào kāichē qù} — Next week I'm going to drive there.
\end{quote}

\begin{cousinweb}
\textbf{\zh{下星期}} (\emph{xià xīngqī}) is built from \zh{下} 'after' and \zh{星期} 'week'. \zh{星} and \zh{期} are inherited from Middle Chinese, \emph{seng} and \emph{gi} (Guangyun, in Baxter's transcription), glossed 'the stars' and 'a meeting; a limit'. The regular sound change from \emph{gi} would give a second tone, qí; why the week says qī is not settled.
\end{cousinweb}

\begin{grammarlens}[title={What you've built}]
\zh{下星期} takes no \zh{个}; \zh{下个月} does. Both go early in the sentence, and both pair with \zh{要} before a verb:

\begin{quote}
\textbf{\zh{下个月我要去中国}} \emph{xià ge yuè wǒ yào qù Zhōngguó} — Next month I'm going to China.
\end{quote}
\end{grammarlens}

\subsection*{Your turn}
\begin{itemize}
\raggedright
  \item \emph{Say it:} \emph{xià xīngqī}
  \item \emph{Say it:} \emph{xià xīngqī}, once more
  \item \emph{Say it:} \zh{下星期我要开车去} — \emph{xià xīngqī wǒ yào kāichē qù} — Next week I'm going to drive there.
  \item \emph{Recall:} say \emph{zhōumò}, then \emph{xià xīngqī}
\end{itemize}

\begin{tcolorbox}[breakable,colback=teal!4,colframe=teal!35!black,title={Before you move on}]
What does \zh{下星期} mean? (\textquotedblleft{}Next week\textquotedblright{}.) Which takes \zh{个}, the week or the month? (\textbf{The month: \zh{下个月}}.)
\end{tcolorbox}
\section[\texorpdfstring{\zh{上个月} (shàng ge yuè)}{上个月 (shàng ge yuè)}]{\zh{上个月} (shàng ge yuè) — last month}

\label{lesson:ZH-C148-shang-ge-yue}



\begin{quote}
Before the new one: last lesson's word, then \zh{了} after a verb.

\begin{itemize}
\raggedright
  \item \emph{Recall:} say \zh{下星期} — \emph{xià xīngqī}, next week — \textbf{R1}, one lesson back
  \item \emph{Recall:} say \zh{我看了书} — \emph{wǒ kàn le shū}, I read the book — \textbf{R4}, recognition at distance
\end{itemize}
\end{quote}

\subsection*{You'll want to know: \zh{上个月}}
\textbf{\zh{上个月}} — \emph{shàng ge yuè} — \textquotedblleft{}last month\textquotedblright{}.

\begin{quote}
\textbf{\zh{上个月我去了中国}} \emph{shàng ge yuè wǒ qù le Zhōngguó} — Last month I went to China.
\end{quote}

\begin{cousinweb}
\textbf{\zh{上个月}} (\emph{shàng ge yuè}) is built from \zh{上} 'above; before', \zh{个} and \zh{月}. \zh{上} is inherited from Middle Chinese \emph{dzyangH} (Guangyun, in Baxter's transcription); the Guangyun glosses this reading 'ruler', so 'above', and 'last' in a row of months, are assumed here, not quoted.
\end{cousinweb}

\begin{grammarlens}[title={What you've built}]
\zh{上} and \zh{下} point through time as through a list read from the top: the month above is the one before, the month below the one after.

\zh{上个月} looks back, so it takes \zh{了} when you say what happened: \zh{上个月}, \zh{我}, \zh{去了}, \zh{中国}.
\end{grammarlens}

\subsection*{Your turn}
\begin{itemize}
\raggedright
  \item \emph{Say it:} \emph{shàng ge yuè}
  \item \emph{Say it:} \emph{shàng ge yuè}, once more
  \item \emph{Say it:} \zh{上个月我去了中国} — \emph{shàng ge yuè wǒ qù le Zhōngguó} — Last month I went to China.
  \item \emph{Say it:} \zh{下个月我要去中国} — \emph{xià ge yuè wǒ yào qù Zhōngguó} — Next month I'm going to China.
  \item \emph{Recall:} say \emph{xià xīngqī}, then \emph{shàng ge yuè}
\end{itemize}

\begin{tcolorbox}[breakable,colback=teal!4,colframe=teal!35!black,title={Before you move on}]
What does \zh{上个月} mean? (\textquotedblleft{}Last month\textquotedblright{}.) And next month? (\textbf{\zh{下个月}}.)
\end{tcolorbox}
\section[dǎsuàn]{dǎsuàn — to plan to; to intend}

\label{lesson:ZH-C148-dasuan}



\begin{quote}
Before the new one: last lesson's word, then a verb this one is built on.

\begin{itemize}
\raggedright
  \item \emph{Recall:} say \zh{上个月} — \emph{shàng ge yuè}, last month — \textbf{R1}, one lesson back
  \item \emph{Recall:} say \emph{suàn}, to work out — \textbf{R2}, a few chapters back
\end{itemize}
\end{quote}

\subsection*{You'll want to know: dǎsuàn}
\emph{dǎsuàn} — \textquotedblleft{}to plan to; to intend\textquotedblright{}.

\begin{quote}
\emph{wǒ dǎsuàn xià ge yuè qù Běijīng} — I'm planning to go to Beijing next month.
\end{quote}

\begin{quote}
\emph{nǐ dǎsuàn zuò shénme} — What are you planning to do?
\end{quote}

Its second character comes later; for now it is a word for the ear.

\begin{cousinweb}
\emph{dǎsuàn} is built from \zh{打}, the \zh{打} of \zh{打电话}, and \emph{suàn} 'to work out', the word you know. Both are inherited from Middle Chinese, \emph{taengX} and \emph{swanX} (Guangyun, in Baxter's transcription), glossed 'to strike' and 'the number of things'. Neither dǎ nor suàn is the sound those readings would regularly give, and how striking and counting came to mean planning is not settled.
\end{cousinweb}

\begin{grammarlens}[title={What you've built}]
\emph{dǎsuàn} comes before a verb, as \zh{要} does: \emph{wǒ dǎsuàn qù}, I plan to go.

\zh{要} says it is going to happen; \emph{dǎsuàn} says it is your plan, and a plan can change.
\end{grammarlens}

\subsection*{Your turn}
\begin{itemize}
\raggedright
  \item \emph{Say it:} \emph{dǎsuàn}
  \item \emph{Say it:} \emph{dǎsuàn}, once more
  \item \emph{Say it:} \emph{nǐ dǎsuàn zuò shénme} — What are you planning to do?
  \item \emph{Say it:} \emph{wǒ dǎsuàn xià ge yuè qù Běijīng} — I'm planning to go to Beijing next month.
  \item \emph{Recall:} say \emph{suàn}, then \emph{dǎsuàn}
\end{itemize}

\begin{tcolorbox}[breakable,colback=teal!4,colframe=teal!35!black,title={Before you move on}]
What does \emph{dǎsuàn} mean? (\textquotedblleft{}To plan to; to intend\textquotedblright{}.) What comes after it? (\textbf{A verb: \emph{dǎsuàn qù}}.)
\end{tcolorbox}
\section[jiānglái]{jiānglái — in the future; some day}

\label{lesson:ZH-C148-jianglai}



\begin{quote}
Before the new one: last lesson's word, then a word about growing up.

\begin{itemize}
\raggedright
  \item \emph{Recall:} say \emph{dǎsuàn}, to plan to — \textbf{R1}, one lesson back
  \item \emph{Recall:} say \zh{长大} — \emph{zhǎngdà}, to grow up — \textbf{R3}, durable
\end{itemize}
\end{quote}

\subsection*{You'll want to know: jiānglái}
\emph{jiānglái} — \textquotedblleft{}in the future; some day\textquotedblright{}.

\begin{quote}
\emph{jiānglái wǒ xiǎng zuò lǎoshī} — Some day I'd like to be a teacher.
\end{quote}

Its first character comes later; for now it is a word for the ear.

\begin{cousinweb}
\emph{jiānglái} is built from jiāng and \zh{来} 'to come'. Both are inherited from Middle Chinese, \emph{tsjang} and \emph{loj} (Guangyun, in Baxter's transcription), glossed 'to send off; to go' and 'to arrive'. 'The time to come' is the pair's own sense, assumed here, not quoted.
\end{cousinweb}

\begin{grammarlens}[title={What you've built}]
\emph{jiānglái} reaches further than \zh{下个月} or \zh{明年}: not a date, but the years ahead.

It goes early, before the person, as the other time words do.
\end{grammarlens}

\subsection*{Your turn}
\begin{itemize}
\raggedright
  \item \emph{Say it:} \emph{jiānglái}
  \item \emph{Say it:} \emph{jiānglái}, once more
  \item \emph{Say it:} \emph{jiānglái wǒ xiǎng zuò lǎoshī} — Some day I'd like to be a teacher.
  \item \emph{Recall:} say \emph{zhǎngdà}, then \emph{jiānglái}
\end{itemize}

\begin{tcolorbox}[breakable,colback=teal!4,colframe=teal!35!black,title={Before you move on}]
What does \emph{jiānglái} mean? (\textquotedblleft{}In the future; some day\textquotedblright{}.) Is it a date? (\textbf{No: the years ahead}.)
\end{tcolorbox}
\section[Practice]{\zh{下个月我要去中国}}

\label{lesson:ZH-C148-practice}



\begin{quote}
The chapter closes here, with two people on a platform, trading plans.

\begin{itemize}
\raggedright
  \item \emph{Recall:} say \emph{jiānglái}, in the future — \textbf{R1}, one lesson back
  \item \emph{Recall:} say \emph{dǎsuàn}, to plan to — \textbf{R1}
  \item \emph{Recall:} say \zh{上个月} — \emph{shàng ge yuè}, last month — \textbf{R2}
\end{itemize}
\end{quote}

\subsection*{The exchange}
\begin{quote}
\textbf{A:} \zh{下星期你要做什么}
\end{quote}

\begin{quote}
\emph{xià xīngqī nǐ yào zuò shénme} — What are you going to do next week?
\end{quote}

\begin{quote}
\textbf{B:} \zh{下星期我要开车去大学}
\end{quote}

\begin{quote}
\emph{xià xīngqī wǒ yào kāichē qù dàxué} — Next week I'm going to drive to the university.
\end{quote}

\begin{quote}
\textbf{A:} \zh{下个月呢}
\end{quote}

\begin{quote}
\emph{xià ge yuè ne} — And next month?
\end{quote}

\begin{quote}
\textbf{B:} \zh{下个月我要去中国}
\end{quote}

\begin{quote}
\emph{xià ge yuè wǒ yào qù Zhōngguó} — Next month I'm going to China.
\end{quote}

\begin{quote}
\textbf{A:} \zh{上个月你也去了中国吗}
\end{quote}

\begin{quote}
\emph{shàng ge yuè nǐ yě qù le Zhōngguó ma} — Did you go to China last month too?
\end{quote}

\begin{quote}
\textbf{B:} \zh{没有} \zh{上个月我在家}
\end{quote}

\begin{quote}
\emph{méiyǒu, shàng ge yuè wǒ zài jiā} — No, last month I was at home.
\end{quote}

\begin{quote}
\textbf{A:} \emph{jiānglái nǐ dǎsuàn zuò shénme} — What do you plan to do in the future?
\end{quote}

\begin{quote}
\textbf{B:} \emph{wǒ xiǎng zuò lǎoshī} — I'd like to be a teacher.
\end{quote}

\zh{要} before each verb says what is coming; \zh{上个月} takes \zh{了} because it is over. \emph{méiyǒu} answers the question about the past, as \zh{没} reports what did not happen.

\begin{tcolorbox}[breakable,skin=enhanced,colback=orange!6,colframe=orange!45!black,boxrule=0.5pt,arc=1mm,left=6pt,right=6pt,top=4pt,bottom=4pt,fonttitle=\bfseries,title={Read this}]
\emph{Read it:} the exchange through once without stopping, then answer the question below
\end{tcolorbox}

\begin{grammarlens}[title={What you've built — next week, next month, some day}]
Time words go first; \zh{要} or \emph{dǎsuàn} goes before the verb:

\begin{itemize}
\raggedright
  \item \textbf{\zh{下星期我要开车去}} \emph{xià xīngqī wǒ yào kāichē qù} — next week I'm going to drive.
  \item \textbf{\zh{下个月我要去中国}} \emph{xià ge yuè wǒ yào qù Zhōngguó} — next month I'm going to China.
  \item \emph{jiānglái wǒ xiǎng zuò lǎoshī} — some day I'd like to teach.
\end{itemize}

\textbf{Listen.} \emph{Xià ge yuè nǐ yào qù nǎr?}

\textbf{Speak.} Say what you are going to do next week, then next month.

\emph{Write it:} \zh{下个月} — \zh{月} is the moon piece of \zh{明}, standing alone
\end{grammarlens}

\subsection*{Your turn}
\begin{itemize}
\raggedright
  \item \emph{Say it:} \zh{下星期你要做什么} — \emph{xià xīngqī nǐ yào zuò shénme} — What are you going to do next week?
  \item \emph{Say it:} the whole exchange, both parts
  \item \emph{Write it:} \zh{下个月我要去中国}
\end{itemize}

\begin{tcolorbox}[breakable,colback=teal!4,colframe=teal!35!black,title={Before you move on}]
Say next month. (\textbf{\zh{下个月}}.) Say you are going to China. (\textbf{\zh{我要去中国}}.) Say the train is about to leave. (\textbf{\zh{火车要开了}}.) Say what you plan to do. (\emph{wǒ dǎsuàn …}.)
\end{tcolorbox}
